\subsection{Native scheme-valued points and genuine section evaluation}
\label{subsec:native-scheme-valued-points}

This point-space construction uses Mathlib schemes, their actual stalk residue
fields, and genuine local sections. No analytic structure sheaf or universal
analytification property is asserted.

\begin{definition}[Native valued points]
\label{scheme:valued-points}
\lean{EnvelopeProperty.Berkovich.SchemeValuedPoint,
EnvelopeProperty.Berkovich.SchemeValuedPoint.eval,
EnvelopeProperty.Berkovich.SchemeValuedPoint.eval_mul,
EnvelopeProperty.Berkovich.SchemeValuedPoint.eval_eq_zero_iff,
EnvelopeProperty.Berkovich.SchemeValuedPoint.eval_restrict}
\leanok
\uses{spectrum:valued-point-equivalence}
For a scheme $X$, a valued point consists of $x\in X$ and a real multiplicative
ultrametric valuation $\nu$ on the actual stalk residue field $\kappa_X(x)$.
For $x\in U$ and a genuine section $s\in\mathcal O_X(U)$, its evaluation is
$|s|=\nu(\operatorname{ev}_{U,x}(s))$.
\end{definition}
\begin{proof}
Use the scheme's residue field and its germ-followed-by-residue evaluation ring
map. Valuation multiplicativity and the ultrametric inequality apply to this
actual ring map. Composing section restriction with a germ gives the same
germ, so genuine section restriction preserves this evaluated value. Its value is zero precisely when the residue evaluation is
zero, equivalently when $x$ is outside the section's scheme basic open.
\end{proof}

\begin{theorem}[Native pullback of valued points]
\label{scheme:valued-point-pullback}
\lean{EnvelopeProperty.Berkovich.SchemeValuedPoint.map,
EnvelopeProperty.Berkovich.SchemeValuedPoint.map_point,
EnvelopeProperty.Berkovich.SchemeValuedPoint.eval_map,
EnvelopeProperty.Berkovich.SchemeValuedPoint.evalTop_map,
EnvelopeProperty.Berkovich.SchemeValuedPoint.map_id,
EnvelopeProperty.Berkovich.SchemeValuedPoint.map_comp}
\leanok
\uses{scheme:valued-points}
A scheme morphism $f:X\to Y$ maps the valued point $(x,\nu)$ to $f(x)$,
with valuation pulled back along the actual residue-field map
$\kappa_Y(f(x))\to\kappa_X(x)$. Evaluating a section of $Y$ at the image
point equals evaluating its genuine section pullback at the original point.
The construction obeys identity and composition on whole valued points.
\end{theorem}
\begin{proof}
Pull back the valuation along the field map induced by the actual local stalk
homomorphism. Naturality of the scheme evaluation map sends every local-section
residue to the residue of its section pullback, so application of $\nu$ proves
the evaluation identity. Actual residue-field identity and composition laws
then prove equality of the whole valuations as well as their underlying points.
\end{proof}

\begin{theorem}[Actual affine residue-field comparison]
\label{scheme:affine-valued-point-comparison}
\lean{EnvelopeProperty.Berkovich.SchemeValuedPoint.toValuedPoint,
EnvelopeProperty.Berkovich.SchemeValuedPoint.ofValuedPoint,
EnvelopeProperty.Berkovich.SchemeValuedPoint.toValuedPoint_prime,
EnvelopeProperty.Berkovich.SchemeValuedPoint.ofValuedPoint_point,
EnvelopeProperty.Berkovich.SchemeValuedPoint.toValuedPoint_ofValuedPoint,
EnvelopeProperty.Berkovich.SchemeValuedPoint.ofValuedPoint_toValuedPoint,
EnvelopeProperty.Berkovich.SchemeValuedPoint.valuedPointEquiv,
EnvelopeProperty.Berkovich.SchemeValuedPoint.toValuedPoint_eval}
\leanok
\uses{scheme:valued-points,spectrum:valued-point-equivalence}
For every commutative ring $A$, native valued points of the actual scheme
$\operatorname{Spec}(A)$ correspond bijectively to $\operatorname{ValuedPoint}(A)$.
The correspondence uses the actual scheme residue-field isomorphism with the
prime's residue field, and preserves evaluation of all original-ring elements
viewed as actual global sections under $\Gamma(\operatorname{Spec}A)\simeq A$.
\end{theorem}
\begin{proof}
The actual scheme residue-field isomorphism identifies the field at a prime
with the quotient fraction field. Pull a valuation back along the inverse
isomorphism to obtain the algebraic valued point; pull it back along the forward
isomorphism for the inverse construction. Inverse isomorphism identities prove
both whole-point inverse laws. The algebra-map/residue-field-isomorphism
compatibility identifies the original-ring element with the residue of its
actual global section, proving the evaluation identity.
\end{proof}

\begin{theorem}[Global evaluation topology and continuous native maps]
\label{scheme:valued-point-topology}
\lean{EnvelopeProperty.Berkovich.SchemeValuedPoint.sectionOpen,
EnvelopeProperty.Berkovich.SchemeValuedPoint.topologyGenerators,
EnvelopeProperty.Berkovich.SchemeValuedPoint.isOpen_point_preimage,
EnvelopeProperty.Berkovich.SchemeValuedPoint.isOpen_sectionOpen,
EnvelopeProperty.Berkovich.SchemeValuedPoint.continuous_point,
EnvelopeProperty.Berkovich.SchemeValuedPoint.continuous_eval,
EnvelopeProperty.Berkovich.SchemeValuedPoint.continuous_map,
EnvelopeProperty.Berkovich.schemeValuedPointFunctor,
EnvelopeProperty.Berkovich.schemeValuedPointProjection}
\leanok
\uses{scheme:valued-points,scheme:valued-point-pullback}
Give native valued points the topology generated by inverse images of scheme
opens under the point projection and by local-section evaluation inverse images
of open subsets of $\mathbb R_{\ge0}$ within their domains. Projection to the
scheme is continuous, evaluation on each open domain is continuous, and every
scheme morphism induces a continuous map on native valued points. These
actual maps define a functor from schemes to topological spaces, with point
projection a natural continuous transformation to the underlying scheme space.
\end{theorem}
\begin{proof}
The topology is generated by the displayed sets, so inverse images of support
opens are open. A local evaluation inverse image on its domain is one of the
generators; taking the subspace proves continuity of the local evaluation.
For a scheme map, inverse images of support generators are generators from its
actual open preimage. Evaluation naturality turns the inverse image of a local
evaluation generator into the generator for the actual pulled-back section.
Thus inverse images of every generator are open, proving continuity. The
whole-point identity and composition laws define the actual topological functor;
its point projection commutes with each underlying scheme map by construction.
\end{proof}

\begin{theorem}[Native affine evaluation homeomorphism]
\label{scheme:affine-valued-point-homeomorphism}
\lean{EnvelopeProperty.Berkovich.SchemeValuedPoint.continuous_toValuedPoint,
EnvelopeProperty.Berkovich.SchemeValuedPoint.ofValuedPoint_eval_algebraMap,
EnvelopeProperty.Berkovich.SchemeValuedPoint.ofValuedPoint_eval_const,
EnvelopeProperty.Berkovich.SchemeValuedPoint.ofValuedPoint_eval_of_const,
EnvelopeProperty.Berkovich.SchemeValuedPoint.continuous_ofValuedPoint,
EnvelopeProperty.Berkovich.SchemeValuedPoint.valuedPointHomeomorph,
EnvelopeProperty.Berkovich.SchemeValuedPoint.spectrumHomeomorph,
EnvelopeProperty.Berkovich.SchemeValuedPoint.toValuedPoint_map,
EnvelopeProperty.Berkovich.SchemeValuedPoint.spectrumHomeomorph_map}
\leanok
\uses{scheme:affine-valued-point-comparison,scheme:valued-point-topology,
spectrum:valued-point-topology,scheme:valued-point-pullback}
The actual affine valued-point equivalence is a homeomorphism between the
native local-section topology and the original-ring evaluation topology.
Composition gives an actual homeomorphism with the seminorm spectrum. Both
comparisons commute with actual $\operatorname{Spec}(\varphi)$ for ring maps.
\end{theorem}
\begin{proof}
Forward continuity follows because every original-ring value is the evaluation
of its actual global section. For reverse continuity, the prime projection is
already continuous. At any prime in a section's domain, the actual affine
structure sheaf gives a basic neighborhood on which the section is a fixed
fraction $a/f$, with $f$ outside every prime in that neighborhood. Restricting
the original global sections commutes with their genuine residue evaluation.
The equality $f(a/f)=a$ in that section ring thus gives the value of the section
as $|a|/|f|$. This is a continuous function of original-ring evaluations on
that nonzero-denominator neighborhood. At a point where the section value
belongs to a specified open interval, this local ratio gives a neighborhood
contained in the required evaluation inverse image. Hence all native topology
generators have open inverse image. For an actual Spec map, scheme global
evaluation naturality and the global-section identification send the section
of $a$ to the section of $\varphi(a)$. Thus every original-ring evaluation
agrees with ring-map pullback. The algebraic valued-point evaluation
injectivity proves equality of whole points and the naturality of both
comparisons. No equality of analytic structure sheaves
or global gluing identification is used.
\end{proof}
